% Responsable: P1
\section{Diseño de la gramática}
\label{sec:gramatica}

\subsection{Producciones}

Partimos de la gramática de ejemplo de Luque~\cite{luque}, que describe oraciones como
\emph{Él mató un elefante en sus pijamas}, y le agregamos lo que pide el enunciado: estructura
sujeto-verbo-objeto, adjetivos, adverbios, frases preposicionales y dos tiempos verbales. La
gramática es $G = (V, \Sigma, R, S)$ con
$V = \{S, NP, Nom, AdjP, VP, PP\}$, terminales en las categorías léxicas
$\Sigma = \{\textit{Prn}, \textit{Det}, \textit{Pos}, \textit{Noun}, \textit{Adj},
\textit{Intens}, \textit{Adv}, \textit{Verb}, \textit{Prep}\}$ y las producciones
\begin{align*}
  S    &\to NP\ VP \\
  NP   &\to \textit{Prn} \mid \textit{Det}\ Nom \mid \textit{Pos}\ Nom \mid NP\ PP \\
  Nom  &\to \textit{Noun} \mid \textit{Noun}\ AdjP                      && \star \\
  AdjP &\to \textit{Adj} \mid \textit{Intens}\ \textit{Adj}             && \star \\
  VP   &\to \textit{Verb}\ NP \mid VP\ PP \mid \textit{Verb}\ \star \mid VP\ \textit{Adv}\ \star \\
  PP   &\to \textit{Prep}\ NP
\end{align*}
Las reglas marcadas con $\star$ son nuevas; las demás vienen de Luque. El único cambio a sus
reglas es que \textit{Det Noun} y \textit{Pos Noun} pasan a ser \textit{Det Nom} y \textit{Pos
Nom}, para que el adjetivo entre sin duplicar reglas. Las categorías léxicas se expanden en
las palabras del léxico (sección~\ref{sec:lexer_concordancia_arboles}).

\begin{table}[h]
  \centering
  \small
  \begin{tabular}{ll}
    \toprule
    Requisito del enunciado & Regla o recurso \\
    \midrule
    Sujeto-verbo-objeto & $S \to NP\ VP$, $VP \to \textit{Verb}\ NP$ \\
    Adjetivos & $Nom \to \textit{Noun}\ AdjP$, $AdjP \to \textit{Adj}$ \\
    Adverbios & $VP \to VP\ \textit{Adv}$, $AdjP \to \textit{Intens}\ \textit{Adj}$ \\
    Frases preposicionales & $NP \to NP\ PP$, $VP \to VP\ PP$, $PP \to \textit{Prep}\ NP$ \\
    Dos tiempos verbales & presente y pretérito, como rasgo \texttt{tiempo} del léxico \\
    Concordancia & rasgos de género, número y persona, verificados aparte \\
    \bottomrule
  \end{tabular}
  \caption{Cómo cubre la gramática cada requisito.}
  \label{tab:cobertura_gramatica}
\end{table}

\subsection{Versión para un parser descendente}

Las reglas $NP \to NP\ PP$, $VP \to VP\ PP$ y $VP \to VP\ \textit{Adv}$ son recursivas por la
izquierda: un parser descendente recursivo que intente expandir $NP$ empezando por $NP$ entra en
un ciclo infinito. Las eliminamos con la transformación estándar
$A \to A\alpha \mid \beta \;\Rightarrow\; A \to \beta A',\ A' \to \alpha A' \mid \varepsilon$,
que genera el mismo lenguaje~\cite{hmu}:
\begin{align*}
  S      &\to NP\ VP                       & NP'    &\to PP\ NP' \mid \varepsilon \\
  NP     &\to NPbase\ NP'                  & AdjOpc &\to AdjP \mid \varepsilon \\
  NPbase &\to \textit{Prn} \mid \textit{Det}\ Nom \mid \textit{Pos}\ Nom
                                           & ObjOpc &\to NP \mid \varepsilon \\
  Nom    &\to \textit{Noun}\ AdjOpc        & VP'    &\to PP\ VP' \mid \textit{Adv}\ VP' \mid \varepsilon \\
  AdjP   &\to \textit{Adj} \mid \textit{Intens}\ \textit{Adj}
                                           & PP     &\to \textit{Prep}\ NP \\
  VP     &\to \textit{Verb}\ ObjOpc\ VP'
\end{align*}
El parser decide qué regla usar mirando solo la palabra actual, con los conjuntos FIRST:
$NP$ y $NPbase$ empiezan con \textit{Prn}, \textit{Det} o \textit{Pos}; $Nom$ con
\textit{Noun}; $AdjP$ con \textit{Adj} o \textit{Intens}; $VP$ con \textit{Verb}; y $PP$ con
\textit{Prep}. Los no terminales auxiliares ($NP'$, $AdjOpc$, $ObjOpc$, $VP'$) no aparecen en
los árboles: sus hijos cuelgan directamente del nodo padre, así que un $VP$ queda como
\textit{Verb} seguido de un $NP$ opcional y de cualquier número de $PP$ y \textit{Adv}. Ese
formato de árbol es el que comparten el parser, el LLM y las métricas.

\subsection{Decisiones de diseño}

\paragraph{\textit{Intens} y \textit{Adv} son categorías distintas.} Si \emph{muy} fuera un
adverbio, en \emph{comió la manzana muy rápidamente} el parser no sabría, al ver un adverbio
después del sustantivo, si empieza un $AdjP$ o un modificador del verbo: \textit{Adv} estaría en
FIRST($AdjOpc$) y en FOLLOW($AdjOpc$), y la gramática dejaría de ser LL(1) en ese punto. Con dos
categorías el conflicto desaparece.

\paragraph{La concordancia no se escribe en la gramática.} Con un número finito de rasgos la
concordancia sí se puede expresar con una CFG, pero a un costo enorme: habría que duplicar $NP$,
$Nom$, $AdjP$ y las categorías \textit{Det}, \textit{Noun} y \textit{Adj} por cada combinación de
género y número (cuatro), y separar los verbos por persona y número (seis combinaciones). Por eso
la gramática acepta \emph{*La elefante grises duerme} y un analizador semántico la rechaza
después, recorriendo el árbol. Es la misma división que en la Fase~1, donde la regla de
\emph{declarar antes de usar} se verificaba con la tabla de símbolos y no con la gramática.

\paragraph{Una ambigüedad que no se puede eliminar.} En \emph{Él mató un elefante en sus
pijamas}, al llegar a \emph{en} el parser puede seguir dentro del objeto (el elefante estaba en
sus pijamas) o cerrarlo y unir la frase preposicional al verbo (él estaba en pijamas).
\textit{Prep} está a la vez en FIRST($NP'$) y en FOLLOW($NP'$), así que la gramática no es
LL(1), y ninguna transformación lo arregla porque la ambigüedad es del español, no de la
gramática: la oración tiene dos árboles (figura~\ref{fig:arboles_pijamas}). Decidimos que el
parser una la frase al sustantivo más cercano, siga devolviendo un solo árbol y lo documentamos
como limitación. La sección~\ref{sec:resultados} mide su costo.

\begin{figure}[h]
  \centering
  \resizebox{0.49\linewidth}{!}{%
  \begin{forest} sintaxis
    [S [NP [NPbase [Prn [él]]]]
       [VP [Verb [mató]]
           [NP [NPbase [Det [un]] [Nom [Noun [elefante]]]]
               [PP [Prep [en]] [NP [NPbase [Pos [sus]] [Nom [Noun [pijamas]]]]]]]]]
  \end{forest}}%
  \hfill
  \resizebox{0.49\linewidth}{!}{%
  \begin{forest} sintaxis
    [S [NP [NPbase [Prn [él]]]]
       [VP [Verb [mató]]
           [NP [NPbase [Det [un]] [Nom [Noun [elefante]]]]]
           [PP [Prep [en]] [NP [NPbase [Pos [sus]] [Nom [Noun [pijamas]]]]]]]]
  \end{forest}}
  \caption{Los dos árboles de \emph{Él mató un elefante en sus pijamas}. Izquierda: la frase
    preposicional se une al objeto (el árbol que construye nuestro parser). Derecha: se une al
    verbo.}
  \label{fig:arboles_pijamas}
\end{figure}

\subsection{Lo que la gramática no puede describir}

La gramática genera un fragmento pequeño y bien definido del español, y el conjunto de prueba
incluye a propósito construcciones correctas que quedan fuera de él
(sección~\ref{sec:set_de_prueba}). Cada una falla por una razón distinta:
\begin{itemize}
  \item \textbf{Sujeto omitido} (\emph{Come la manzana}, \emph{Leímos un libro}). La regla
    $S \to NP\ VP$ exige un sujeto explícito. Agregar $S \to VP$ es fácil, pero obliga a
    deducir el sujeto de la persona del verbo.
  \item \textbf{Orden libre} (\emph{Leen las niñas un libro}). La gramática fija el orden
    sujeto-verbo-objeto. Cada orden adicional necesita sus propias reglas, y decidir cuál
    sintagma es el sujeto ya no depende de la posición sino de la concordancia.
  \item \textbf{Coordinación} (\emph{El perro y el gato duermen}). Faltan reglas como
    $NP \to NP\ \textit{Conj}\ NP$, que además cambian la concordancia: dos sustantivos
    singulares forman un sujeto plural.
  \item \textbf{Subordinadas} (\emph{La niña que lee un libro come}). Necesitan oraciones dentro
    de sintagmas nominales ($NP \to NP\ \textit{que}\ S/NP$) con un hueco donde falta el sujeto o
    el objeto. Las CFG pueden describir algunos de estos casos, pero la gramática crece mucho y
    se vuelve muy ambigua.
  \item \textbf{Vocabulario cerrado} (\emph{El perro corre en el parque}). Basta una palabra que
    no esté en el léxico para que la oración se rechace, aunque su estructura esté en la
    gramática.
\end{itemize}
Más allá de este fragmento, hay construcciones de algunas lenguas, como las dependencias
cruzadas del suizo alemán, que ninguna CFG puede generar~\cite{shieber}. Nuestro fragmento no
llega a ese límite: lo que lo restringe es la cobertura, no el poder expresivo de las CFG.
